\documentclass[10pt]{article}

\usepackage[utf8]{inputenc}

\usepackage[T1]{fontenc}

\usepackage{amsmath}

\usepackage{amsfonts}

\usepackage{amssymb}

\usepackage[version=4]{mhchem}

\usepackage{stmaryrd}

\usepackage{bbold}

\usepackage{mathrsfs}



\begin{document}

ренцируемую функцию $\psi(\boldsymbol{r}, t)$, удовлетворяюшую В. у. в полупространстве $t>0$ и начальным условиям



\[

\left.\psi\right|_{t=0}=\varphi_{1}(\boldsymbol{r}),\left.\quad \frac{\partial \psi}{\partial t}\right|_{t=0}=\varphi_{2}(\boldsymbol{r}),

\]



где $\varphi_{1}(\boldsymbol{r})$ и $\varphi_{2}(\boldsymbol{r})$ - заданные функции. Классич. решение дается Кирхгофа формулой $(n=3)$, Пуассона формулой $(n=2)$ или Д'Аламбера формулой $(n=1)$. Рассматривают также смешанную задачу, описывающую колебания ограниченного объема $V$.



При перечислении нелинейных обобщений В. у. необходимо проявлять нек-рую сдержанность, с тем чтобы при этом не утрачивалась связь с исходным В. у. В этом смысле единственным терминологически точным обобщением является внесение зависимости скорости $c$ от волновой функции в уравнение (1) или (5). Однако часто к нелинейным В. у. относят любые уравнения, вырождающиеся в линейные В. у. при устранении нелинейности или линеаризации. Наиболее известны нелинейное уравнение Клейна - Гордона



\[

\square \psi=m^{2} \psi+F(\psi)

\]



обобщающее линейное Клейна - Гордона уравнение, и нелинейное уравнение Гельмгольца



\[

\Delta \psi+k^{2} \psi=F\left(|\psi|^{2}\right) \psi,

\]



учитывающее зависимость волнового числа от квадрата волновой функции.



Нелинейные В.у. позволяют описать взаимодействие волн, возникновение и эволюцию ударных волн и солитонов и т. Д.



Лит.: [1] Физическая энциклопедия, т. 1, М., 1988, с. 312-13; [2] М о рс Ф., Феш бах Г., Методы теоретической физики, пер. с англ., т. 1-2, М., 1958-60; [3] Влад им и ров В. С., Уравнения математической физики, 4 изд., М., 1981; [4] У и зе м Дж., Линейные и нелинейные волны, пер. с англ., М., 1977.



М. А. Миллер, Е. И. Якубович.



ВОЛНОВОЕ УРАВНЕНИЕ; ас имптотические ре ш е н и я - решения волнового уравнения



\[

\square_{c} \psi=0, \square_{c}=\partial_{t}^{2}-c^{2}(x) \Delta, 0<c(x) \in C^{\infty}\left(\mathbb{R}_{x}^{n}\right),

\]



описывающие эволюцию коротковолновых пакетов, которые удовлетворяют начальным данным



\[

\begin{gathered}

\left.\psi\right|_{t=0}=\varphi_{1} \exp \left(i S_{0} / h\right),\left.\quad \psi_{t}\right|_{t=0}=i h^{-1} \varphi_{2} \exp \left(i S_{0} / h\right), \\

\varphi_{1,2} \in C_{0}^{\infty}\left(\mathbb{R}_{x}^{n}\right), \quad S_{0} \in C^{\infty}\left(\mathbb{R}_{x}^{n}\right),

\end{gathered}

\]



где $h-$ малый параметр, $\nabla S_{0} \neq 0$ при $x \in \operatorname{supp} \varphi_{1} \cup \operatorname{supp} \varphi_{2}$. Асимптотич. решение такой задачи Коши (то есть функция, удовлетворяющая уравнению и начальным условиям с точностью $o(1))$ имеет вид



\[

\psi_{a}=\frac{1}{2} \sum_{+,-} K^{ \pm}\left[\frac{\varphi_{1}(\alpha) \mp \varphi_{2}(\alpha)}{\left|\nabla S_{0}(\alpha)\right|} \sqrt{\frac{c(\alpha)}{c\left(X^{ \pm}(\alpha, t)\right)}}\right] \exp \left(i \delta^{ \pm} / h\right),

\]



где $K^{ \pm}$- канонич. оператор Маслова на лагранжевом многообразии $L_{t}^{ \pm}=g^{ \pm}(t) L_{0}$, полученном сдвигом многообразия



\[

L_{0}=\left\{x=\alpha \in \mathbb{R}^{n}, p=\nabla S_{0}(\alpha)\right\} \subset \mathbb{R}_{x}^{n} \times \mathbb{R}_{p}^{n}

\]



вдоль траекторий $X^{ \pm}(\alpha, t), P^{ \pm}(\alpha, t)$ гамильтоновых систем $\dot{x}=H_{p}^{ \pm}, \dot{p}=-H_{x}^{ \pm}$, с отмеченной точкой $\left\{x=X^{ \pm}\left(\alpha_{0}, t\right)\right.$, $\left.p=P^{ \pm}\left(\alpha_{0}, t\right)\right\}, \alpha_{0}$ - нек-рое фиксированное значение параметра $\alpha, H^{ \pm}(x, p)= \pm c(x)|p|, g^{ \pm}(t)-$ фазовые потоки гамильтоновых систем с этими гамильтонианами,



\[

\delta^{ \pm}=\int_{0}^{t}\left(P^{ \pm} X^{ \pm}-H^{ \pm}\right) d \tau+\pi h \operatorname{Ind} l / 2 .

\]



В последней формуле интегрирование ведется вдоль траектории



\[

l=\left\{x=X^{ \pm}\left(\alpha_{0}, \tau\right), p=P^{ \pm}\left(\alpha_{0}, \tau\right), 0 \leqslant \tau \leqslant t\right\},

\]



Ind $l$ - индекс Маслова этой траектории. С помощью несколько более громоздких формул вычисляются поправки к асимптотич. решению с любой степенью точности по $h$. Ограничение $\nabla S_{0} \neq 0$ связано с тем, что гамильтонианы $H^{ \pm}(x, p)$ совпадают и являются негладкими при $p=0$. Интегрирование начальных данных по $1 / h$ переводит их в функции $2 \pi \varphi_{1} \delta\left(S_{0}\right)$ и $2 \pi \varphi_{2} \delta^{\prime}\left(S_{0}\right)$, а функцию $O\left(h^{N}\right), N \gg 1,-$ в гладкую функцию. Указанные выше асимптотики тесно связаны с построением асимптотич. решения по гладкости, то есть функции, к-рая отличается от точного решения на достаточное число раз непрерывно дифференцируемую функцию. Для начальных данных, носитель к-рых принадлежит нек-рому компакту $\mathscr{K} \subset \mathbb{R}_{x}^{n}$, это делается с помошью параметрикса задачи Коши



\[

\square_{c} \psi=0,\left.\quad \psi\right|_{t=0}=\psi_{0} \in H_{s},\left.\psi_{t}\right|_{t=0}=0,

\]



где $H_{s}\left(\mathbb{R}_{x}^{n}\right)$ - пространство Соболева, $s \in \mathbb{R}$. П а р а м е т р и к сом задачи называется последовательность операторов $T_{N}\left(x,-i \frac{\partial}{\partial x}, t\right)$ такая, что композиция $M_{N}=\square_{c} T_{N}$ является сглаживаюшим оператором в шкале $H_{S}$ :



\[

\| M_{N}\left(x,-i \frac{\partial}{\partial x}, \Psi \|_{H_{k+s}} \leqslant \text { const } \| t\right) \psi \|_{H_{s}},

\]



$k=k(N) \rightarrow \infty$ при $N \rightarrow \infty$, причем



\[

\left.T_{N}\right|_{t=0}=\mathbb{P}(x), \quad \partial T_{N} /\left.\partial t\right|_{t=0}=0,

\]



где $\mathscr{P} \in C_{0}^{\infty}\left(\mathbb{R}_{x}^{n}\right)$ и $\mathscr{y}=1$ при $x \in \mathscr{K}$. При малых $t$ символ $T_{N}(x, \xi, t)$ имеет вид



\[

\begin{gathered}

T_{N}=\sum_{+,-} \exp \left(i|\xi| S^{ \pm}(x, t, \omega)-i \xi x\right) \times \\

\times\left[a_{0}^{ \pm}+\chi(|\xi|) a_{1}^{ \pm}|\xi|^{-1}+\ldots+\chi(|\xi|) a_{N}^{ \pm}|\xi|^{-N}\right],

\end{gathered}

\]



где $\chi(\tau) \in C^{\infty}\left(\mathbb{R}^{1}\right), \chi(\tau)=0$ при $|\tau| \leqslant 1 / 2$ и $\chi=1$ при $|\tau| \geqslant 1$, $\omega=\xi|\xi|^{-1}$, функции $S^{ \pm}$удовлетворяют уравнениям Гамильтона-Якоби



\[

S_{t}^{ \pm}+H^{ \pm}\left(x, \nabla S^{ \pm}\right)=0,\left.S^{ \pm}\right|_{t=0}=\omega x,

\]



(гладкие) функции $a_{j}(x, t, \omega)$ суть решения соответствуюших уравнений переноса. Приведенная формула имеет смысл при тех $t$, при к-рых семейства многообразий $L_{t}^{ \pm}(\omega)=g^{ \pm}(t) L_{0}(\omega)$, полученные из многообразия $L_{0}(\omega)=$ $=\left\{x \in \mathbb{R}^{n}, p=\omega \in \mathbb{S}^{n-1}\right\}$, гладко проектируются на $\mathbb{R}_{x}^{n}$; при произвольных конечных временах символ $T_{N}$ дается канонич. оператором Маслова на этих многообразиях, отвечаюшим параметру $t$. Асимптотика решения указанной задачи Коши по гладкости (с точностью до функций из $H_{s+N}$ ) имеет вид $\psi \sim T_{N} \psi_{0}$. В частности, если $\psi_{0}=\varphi_{0} \theta\left(S_{0}\right)$, где $\varphi_{0} \in C_{0}^{\infty}\left(\mathbb{R}_{x}^{n}\right), S_{0} \in C^{\infty}\left(\mathbb{R}_{x}^{n}\right), \nabla S_{0} \neq 0$ при $S_{0}=0, \theta-$ функция Хевисайда, то интеграл $T_{N} \psi_{0}$, описывающий эволюцию разрыва, с точностью до непрерывных функций по методу стационарной фазы при малых $t$ приводится к классич. виду



\[

\sum_{+,-} \varphi^{ \pm}(x, t) \theta\left(S^{ \pm}(x, t)\right),

\]



где $S^{ \pm}(x, t)$ - решения уравнений Гамильтона-Якоби, $\left.S^{ \pm}\right|_{t=0}=S_{0}, \varphi^{ \pm}-$решения уравнений переноса. Аналогичным образом могут быть изучены особенности фундаментального решения В. у. и описана «метаморфоза разрыва», заключающаяся в изменении вида особенности при прохождении фокальных точек. Параметрикс позволяет также доказать теорему существования и единственности решения. А именно, подстановка



\[

\psi=T_{N} \psi_{0}+\int_{0}^{t} \int_{0}^{t-\tau} T_{N}(x,-i \partial / \partial x, \mu) g(x, \tau) d \mu d \tau

\]



приводит к интегральному уравнению



\[

\int_{0}^{t} \int_{0}^{t-\tau} M_{N}(x,-i \partial / \partial x, \mu) g(x, \tau) d \mu d \tau+g=M_{N} \psi_{0}

\]





\end{document}